\textbf{结论名称：}独立性检验$\chi^2$公式与$2\times2$列联表

\textbf{适用条件：}两个分类变量，样本独立随机，各单元格期望频数均$\ge 5$，总频数$n\ge 40$。

\textbf{变量说明：}设$2\times2$列联表为
\[
\begin{array}{c|cc|c}
 & B_1 & B_2 & \text{合计} \\ \hline
A_1 & a & b & a+b \\
A_2 & c & d & c+d \\ \hline
\text{合计} & a+c & b+d & n
\end{array}
\]
其中$a,b,c,d$为实际观测频数，$n=a+b+c+d$。

\textbf{核心公式：}
\[
\boxed{\chi^2 = \frac{n(ad-bc)^2}{(a+b)(c+d)(a+c)(b+d)}}
\]
或等价地$\displaystyle\chi^2 = \sum \frac{(O-E)^2}{E}$，其中$E$为对应单元格的期望频数（如$E_{A_1B_1}=\frac{(a+b)(a+c)}{n}$），自由度$df=1$。

\textbf{使用场景：}推断两个分类变量是否独立，例如“吸烟与患肺癌是否有关”“性别与某节目喜好是否有关”等。

\textbf{简单例子：}调查200人，数据如下表：
\[
\begin{array}{c|cc|c}
 & \text{患肺癌} & \text{未患} & \text{合计} \\ \hline
\text{吸烟} & 60 & 40 & 100 \\
\text{不吸烟} & 30 & 70 & 100 \\ \hline
\text{合计} & 90 & 110 & 200
\end{array}
\]
计算得$\chi^2 = \frac{200(60\times70-40\times30)^2}{100\times100\times90\times110} \approx 18.18$。若显著性水平$\alpha=0.05$，查表得临界值$\chi^2_{0.05}(1)=3.841$，因为$18.18>3.841$，故拒绝“相互独立”的原假设，可认为吸烟与患肺癌显著相关。